\section{Semilinear localization of affiliated coefficients}
\label{sec:analytic-localization}
\label{coeff:analytic}

Write \(\mathcal N(J)\) for the group von Neumann algebra of a discrete
group \(J\), with its canonical trace, and write \(\mathcal U(J)\) for
its algebra of affiliated operators.  The division fields used here are
embedded in these affiliated algebras. Our target is the following
localization statement for a finite semilinear coefficient system.

\begin{lemma}[Semilinear finite-dimensional support]
\label{coeff:analytic-support}
Let \(G\) be a torsion-free word-hyperbolic group, let
\(h:G\to\Z\) be a homomorphism, and put \(L=\ker h\).
Let \(D\le G\) be quasiconvex and put \(S=D\cap L\).
Choose \(g\in G\) and \(g_D\in D\) such that
\[
 h(g)=h(g_D)\ne0.
\]
Suppose that the division fields in use are embedded in the affiliated
group algebras, and that, for an integer \(r\ge1\), a row
\(\lambda\in D_L^{\,1\times r}\) satisfies
\begin{equation}\label{coeff:analytic-semilinear}
 g\lambda g_D^{-1}=\lambda B,\qquad
 B\in\operatorname{GL}_r(D_S).
\end{equation}
Then
\begin{equation}\label{coeff:analytic-support-conclusion}
 \lambda\in\sum_{xS\in\mathcal P}x\,\mathcal U(S)^{\,1\times r},
 \qquad
 \mathcal P=
 \{xS:\ x\in L,\ g^m\in xDx^{-1}
                    \text{ for some }m\in\Z\setminus\{0\}\}.
\end{equation}
The set \(\mathcal P\) is finite.  In fact, the same assertion holds
for any affiliated row \(\lambda\) and any affiliated invertible
matrix \(B\) satisfying \eqref{coeff:analytic-semilinear}.
\end{lemma}

Here is the normalization behind the proof. Put
\(\alpha=\operatorname{Ad}g_D\), and first consider the case in
which the map \(w\mapsto\lambda w\) on affiliated \(S\)-columns
has zero kernel. We seek a positive invertible affiliated matrix
\(P\) over \(S\) such that
\[
 P=B\alpha(P)B^*.
\]
Then
\[
 Z=\lambda P^{1/2},\qquad
 u=P^{-1/2}B\alpha(P)^{1/2}
\]
satisfy \(uu^*=u^*u=1\) and \(gZg_D^{-1}=Zu\).
Thus changing the positive form makes the coefficient transport
unitary. It follows that \(gZZ^*g^{-1}=ZZ^*\). Hyperbolicity and
the nonzero height of \(g\) will force this invariant affiliated
operator to be scalar, so every entry of \(Z\) is bounded.
These entries generate a closed right \(S\)-module in \(\ell^2(L)\)
of von Neumann dimension at most \(r\). Its invariance under height
translation excludes infinite
coset orbits from the support, and quasiconvexity of \(D\) leaves
only finitely many periodic cosets.

The invariant-form strategy parallels the bounded unitarization
argument of Sz.-Nagy and Dixmier
\cite[Section~5, Theorem~6]{Dixmier1950Means}.
The positive mean constructed below supplies \(P\). Its inputs
need not have a common operator-norm bound, so we average bounded
resolvent expressions and control the spectral tails of the inputs
and their inverses. When the column kernel is nonzero, the proof
first represents its complement by a projection over \(\mathcal N(S)\)
and carries out the same construction between the corresponding
corners.

The affiliated conclusion in \eqref{coeff:analytic-support-conclusion}
is sufficient for Proposition~\ref{close:join-injection}. There the
Bass--Serre homology injection is first established over a subgroup
division field. Extending scalars to its affiliated algebra preserves
that injection and its nonzero classes: choosing a vector-space
basis identifies the scalar extension with a direct sum of copies
of the nonzero affiliated algebra. Injectivity into ambient homology
still requires the separate comparison using actual row chains and
coset independence in that proposition.

\subsection{Measure estimates and a projection inequality}

Let \((\mathcal M,\tau)\) be a finite von Neumann algebra with a faithful
normal finite trace.  The trace need not have total mass one.  On matrices
we use the compatible \emph{unnormalized} traces
\[
 \tau_r((a_{ij}))=\sum_{j=1}^r\tau(a_{jj}).
\]
In particular, for a bounded rectangular matrix \(a\),
\(\tau_r(a^*a)=\tau_k(aa^*)\), with the appropriate source and target
matrix sizes.  Corners carry the restriction of these traces.  Products
and sums of affiliated operators mean their closed affiliated products
and sums.

A family \(\mathcal A\) of affiliated operators is \emph{tight} if
\[
 \lim_{K\longrightarrow\infty}
 \sup_{a\in\mathcal A}\tau\bigl(1_{(K,\infty)}(|a|)\bigr)=0.
\]
For rectangular operators the trace is taken in the source corner;
the nonzero spectral distributions of \(|a|\) and \(|a^*|\) agree
with the compatible traces.  We use convergence in measure: \(a_j\to0\)
if, for every \(\varepsilon,\eta>0\), eventually there is a projection
\(q\) with \(\tau(1-q)<\eta\) and \(\|a_jq\|<\varepsilon\).

\begin{lemma}[Measure calculus]\label{coeff:analytic-calculus}
The affiliated algebra of a finite von Neumann algebra is complete for
convergence in measure.  Products of tight families are tight.  If
\(a_j\to0\) in measure and \(b_j,c_j\) range over tight families, then
\(b_ja_jc_j\to0\) in measure.  These assertions also hold uniformly over
an additional index.

Positive square roots are continuous in measure on positive affiliated
operators, and inversion is continuous in measure on positive affiliated
units.  A uniformly norm-bounded family that tends to zero in measure
tends to zero in \(L^1(\tau)\); the same assertion holds uniformly.
\end{lemma}

\begin{proof}
The complete topological \(*\)-algebra assertion and the singular-number
estimates used below are the standard measurable-operator calculus;
see \cite[Sections~1.4--1.5, Lemma~2.5 and Lemma~3.1]{FackKosaki1986}.
For clarity, the estimates that provide all the uniformity we need are
\[
 \mu_{s+t}(ab)\le \mu_s(a)\mu_t(b),\qquad
 \mu_{s+t}(a+b)\le \mu_s(a)+\mu_t(b),
 \quad s,t>0,
\]
where
\[
 \mu_s(a)=
 \inf\{\|aq\|:\ q\text{ a projection},\ \tau(1-q)\le s\}.
\]
Their proofs cut the two operators to bounded parts and add the
trace codimensions; polar decomposition identifies the trace
codimensions of the source and range cuts.  Tightness says precisely
that \(\sup_{a\in\mathcal A}\mu_s(a)<\infty\) for every \(s>0\).
Convergence in measure says that \(\mu_s(a_j)\to0\) for every \(s>0\).
The two displayed estimates therefore prove the product and uniform
assertions.

We recall how to obtain the stated functional calculus without any
norm bounds on the original operators.  For positive \(a_j\to a\),
the resolvent identity and the product estimate give
\[
 (1+a_j)^{-1}-(1+a)^{-1}
 =(1+a_j)^{-1}(a-a_j)(1+a)^{-1}\longrightarrow0
\]
in measure.  On bounded self-adjoint contractions, polynomial
approximation gives continuity of continuous functional calculus in
measure.  Apply this to the positive contractions
\(q_j=(1+a_j)^{-1}\), with limit \(q=(1+a)^{-1}\).
Here is the endpoint control needed for unbounded scalar functions.
If \(q\) has no spectral atom at \(1\), then for every \(\eta>0\)
one can choose \(b<1\) with \(\mu_{\eta/2}(q)<b\).
The sum inequality gives
\(\mu_\eta(q_j)\le\mu_{\eta/2}(q)+\mu_{\eta/2}(q_j-q)\),
so the spectral mass of \(q_j\) sufficiently near \(1\) is at most
\(\eta\), for all large \(j\).  Applying the same argument to
\(1-q_j\) controls the endpoint \(0\) when \(q\) has no atom there.
The latter condition always holds; the former holds when \(a\) is
a positive unit.  Cut off the functions
\((q^{-1}-1)^{1/2}\), and, in the unit case,
\((q^{-1}-1)^{-1}\), at these endpoints.  The clipped functions are
continuous on \([0,1]\), and the cutoff errors have arbitrarily small
support trace.  This proves square-root and inverse continuity.
The same proof is uniform whenever the required tail bounds are
uniform.

Finally, if \(\|z_j\|\le C\), then
\[
 \|z_j\|_1
 \le \varepsilon\tau(1)
       +C\,\tau\bigl(1_{(\varepsilon,\infty)}(|z_j|)\bigr).
\]
First use convergence in measure and then let \(\varepsilon\) tend to
zero.
\end{proof}

A \emph{positive affiliated unit} in a corner is a positive affiliated
operator whose support is the identity of that corner; its inverse is
allowed to be unbounded.  For two such units define
\begin{equation}\label{coeff:analytic-F}
 F(P,X)=P^{1/2}(P+X)^{-1}P^{1/2}
       =\bigl(1+P^{-1/2}XP^{-1/2}\bigr)^{-1}.
\end{equation}
This is a bounded positive contraction.

\begin{lemma}[Spectral-projection comparison]
\label{coeff:analytic-projection}
Let \(T,X\) be positive affiliated units in \(\mathcal M\), and let
\(t>0\).  For \(p_+=1_{(t,\infty)}(T)\) and
\(p_-=1_{(0,t)}(T)\),
\begin{align}
 \tau\bigl(p_+F(T,X)\bigr)
     &\ge \tau\bigl(p_+F(t1,X)\bigr), \label{coeff:analytic-high}\\
 \tau\bigl(p_-F(T,X)\bigr)
     &\le \tau\bigl(p_-F(t1,X)\bigr). \label{coeff:analytic-low}
\end{align}
\end{lemma}

\begin{proof}
First assume that \(T,X\) and their inverses are bounded.  Put
\(p=p_+\) and \(q=1-p\), and decompose with respect to \(p+q\):
\[
 T=\begin{pmatrix}A&0\\0&B\end{pmatrix},\qquad
 X=\begin{pmatrix}X_{11}&X_{12}\\X_{21}&X_{22}\end{pmatrix}.
\]
Here \(A\ge tp\) and \(B\le tq\).  The positive Schur complements
\[
 K_B=X_{11}-X_{12}(B+X_{22})^{-1}X_{21},\qquad
 K_t=X_{11}-X_{12}(tq+X_{22})^{-1}X_{21}
\]
satisfy \(K_B\le K_t\).  Block inversion and cyclicity of the trace give
\[
 \tau(pF(T,X))=\tau\bigl(A(A+K_B)^{-1}\bigr).
\]
For positive invertible \(A,K\) in \(p\mathcal Mp\), set
\[
 \psi(A,K)=\tau\bigl(A(A+K)^{-1}\bigr).
\]
This function is increasing in \(A\), because
\[
 \psi(A,K)=\tau(p)
  -\tau\bigl(K^{1/2}(A+K)^{-1}K^{1/2}\bigr),
\]
and is decreasing in \(K\), by its defining expression with
\(A^{1/2}\) on the two sides of the inverse.  Consequently
\[
 \psi(A,K_B)\ge\psi(tp,K_B)\ge\psi(tp,K_t)
                 =\tau(pF(t1,X)),
\]
which proves \eqref{coeff:analytic-high}.  For \(p=p_-\), the comparisons
\(A\le tp\), \(B\ge tq\), and \(K_B\ge K_t\) reverse both inequalities,
giving \eqref{coeff:analytic-low}.  The cases \(p=0\) or \(q=0\) follow
from the same argument with the empty block omitted.

For affiliated units, clip \(T\) and \(X\) spectrally into
\([1/R,R]\), with \(1/R<t<R\).  The high and low spectral projections
of the clipped \(T\) at \(t\) are the original projections.  The clipped
operators converge in measure to \(T,X\), and their inverses converge
to the corresponding inverses.  Lemma~\ref{coeff:analytic-calculus}
and \eqref{coeff:analytic-F} show that the resulting positive
contractions converge in \(L^1\).  Both inequalities pass to the limit.
\end{proof}

\subsection{An affiliated positive mean}

Fix a shift-invariant mean \(m\) on \(\ell^\infty(\Z)\).  Such a mean
is obtained as a weak-* cluster point of the averaging functionals
on the intervals \([-N,N]\): translating an interval by a fixed
integer changes its average by a quantity tending to zero.
For a bounded
sequence \(Y_n\in\mathcal M\), define \(\mathfrak M_nY_n\in\mathcal M\)
by duality with the predual:
\[
 \varphi(\mathfrak M_nY_n)=m\bigl(n\mapsto\varphi(Y_n)\bigr),
 \qquad \varphi\in\mathcal M_*.
\]
This mean is positive and unital, commutes with fixed bounded left
and right multiplications, and is shift invariant.  It also satisfies
\begin{equation}\label{coeff:analytic-mean-L1}
 \|\mathfrak M_nY_n\|_1\le\sup_n\|Y_n\|_1.
\end{equation}
Indeed, pair with an arbitrary element of the unit ball of
\(\mathcal M\), using the dual characterization of the \(L^1\)-norm.
No interchange of a nonuniform monotone limit with this mean will be
used.

For finitely many bounded positive invertible inputs and a weighted
average, the equation below defines the induced operator mean
associated with the binary harmonic mean
\cite[Proposition~6.2]{Palfia2013Means}. We prove the affiliated
version needed here, under tightness of the inputs and their inverses
and with the fixed shift-invariant mean \(m\).

\begin{proposition}[Affiliated positive mean]
\label{coeff:analytic-mean}
Let \(X_n\), \(n\in\Z\), be positive affiliated units in a finite von
Neumann algebra.  Suppose that both \(\{X_n\}\) and \(\{X_n^{-1}\}\)
are tight.  There is a unique positive affiliated unit \(P\) satisfying
\begin{equation}\label{coeff:analytic-mean-equation}
 \mathfrak M_nF(P,X_n)=\tfrac12\,1.
\end{equation}
The construction is covariant under congruence by affiliated
invertible operators, including affiliated invertible maps between
corners with their restricted traces.
\end{proposition}

\begin{proof}
We divide the argument into bounded construction, uniform tails,
comparison, and passage to the limit.

For an integer \(R\ge1\), let \(X_{n,R}\) be the spectral clipping of
\(X_n\) into \([R^{-1},R]\).  Define
\[
 H(P,X)=(P^{-1}+X^{-1})^{-1},\qquad
 \Phi_R(P)=2\mathfrak M_nH(P,X_{n,R}).
\]
Inversion reverses positive order, so \(\Phi_R\) is order preserving.
It maps the interval \([R^{-1}1,R1]\) into itself.  Starting from
\(P^{(0)}=R1\), the iterates \(P^{(j+1)}=\Phi_R(P^{(j)})\) decrease and
remain in that interval.  Their strong limit \(P_R\) lies in the same
interval, and normality of the finite trace gives convergence in
\(L^1\).

This convergence passes through \(\Phi_R\) uniformly in \(n\).
In fact,
\[
 H(P,X)=X-X(P+X)^{-1}X,
\]
and the resolvent identity gives, on the fixed interval in question,
\[
 \|H(P,X)-H(Q,X)\|_1\le C_R\|P-Q\|_1
\]
with a constant independent of \(X=X_{n,R}\).
Equation~\eqref{coeff:analytic-mean-L1} therefore shows that
\(P_R=\Phi_R(P_R)\).  The identity
\begin{equation}\label{coeff:analytic-parallel-sum}
 H(P,X)=P^{1/2}\bigl(1-F(P,X)\bigr)P^{1/2}
\end{equation}
and bounded invertibility of \(P_R\) imply
\begin{equation}\label{coeff:analytic-clipped-equation}
 \mathfrak M_nF(P_R,X_{n,R})=\tfrac12\,1.
\end{equation}

The clipped families \(\{X_{n,R}:n\in\Z,R\ge1\}\) and their inverses
are tight: above any threshold \(K>1\), clipping can only decrease
the relevant spectral tail.  Consequently
\begin{align*}
 \sup_{n,R}\|1-F(t1,X_{n,R})\|_1&\longrightarrow0
                                      &&(t\longrightarrow\infty),\\
 \sup_{n,R}\|F(s1,X_{n,R})\|_1&\longrightarrow0
                                      &&(s\longrightarrow0).
\end{align*}
For example, cutting \(X_{n,R}\) at a fixed \(K\) bounds the first
quantity by
\[
 \frac{K}{t}\tau(1)
   +\sup_{n,R}\tau(1_{(K,\infty)}(X_{n,R})).
\]
For the second, cut \(X_{n,R}^{-1}\) at \(K\); the analogous bound is
\(sK\tau(1)\) plus its tail trace.

Apply Lemma~\ref{coeff:analytic-projection} to
\(p=1_{(t,\infty)}(P_R)\) and use
\eqref{coeff:analytic-clipped-equation}.  It gives
\[
 \tfrac12\tau(p)
 \ge \tau(p)-\sup_n\|1-F(t1,X_{n,R})\|_1.
\]
For \(q=1_{(0,s)}(P_R)\), the low-projection inequality gives
\[
 \tfrac12\tau(q)
 \le \sup_n\|F(s1,X_{n,R})\|_1.
\]
Thus both \(\{P_R\}\) and \(\{P_R^{-1}\}\) are tight, uniformly in \(R\).

We next remove clipping from the equation.  Put
\[
 q_{n,R}=1_{(0,R^{-1})\cup(R,\infty)}(X_n).
\]
Two-sided tightness implies \(\sup_n\tau(q_{n,R})\to0\).
The resolvent identity expresses
\begin{align*}
 F(P_R,X_n)-F(P_R,X_{n,R})
 ={}&P_R^{1/2}(P_R+X_n)^{-1}\\
 &{}\cdot(X_{n,R}-X_n)(P_R+X_{n,R})^{-1}P_R^{1/2}.
\end{align*}
The middle factor is supported on \(q_{n,R}\).  Polar decomposition
shows that the support trace of a product cannot exceed the support
trace of any one of its factors.  The displayed difference, being a
difference of positive contractions, has norm at most two.  Hence
\[
 \sup_n\|F(P_R,X_n)-F(P_R,X_{n,R})\|_1
       \le2\sup_n\tau(q_{n,R})\longrightarrow0.
\]
It follows that
\begin{equation}\label{coeff:analytic-approx-error}
 E_R:=\mathfrak M_nF(P_R,X_n)-\tfrac12\,1
       \quad\hbox{satisfies}\quad \|E_R\|_1\longrightarrow0.
\end{equation}

We record precisely the congruence rule.  If \(C\) is an affiliated
invertible operator, set \(P'=CPC^*\) and \(X'=CXC^*\).  Then
\begin{equation}\label{coeff:analytic-congruence}
 F(P',X')=U F(P,X)U^*,\qquad
 U=(P')^{-1/2}CP^{1/2}.
\end{equation}
The affiliated algebra identities \(UU^*=U^*U=1\) show that \(U\)
is a bounded unitary.  The formula follows by taking the inverse of
\(C(P+X)C^*\).  For maps between corners the same calculation makes
\(U\) a partial unitary between their identities.  Crucially, \(U\)
depends on \(P,C\) but not on \(X\); errors in a mean equation are
therefore conjugated by one fixed unitary and keep their \(L^1\)-norm.

We now prove that \(P_R\) is Cauchy in measure, rather than appealing
to any compactness assertion for tight families.  Given \(R,S\), put
\[
 C=P_R^{-1/2},\qquad
 T=C P_S C,\qquad X'_n=C X_n C.
\]
The equations with errors \(E_R,E_S\), after this congruence, become
\[
 \mathfrak M_nF(1,X'_n)=\tfrac12\,1+E_R,\qquad
 \mathfrak M_nF(T,X'_n)=\tfrac12\,1+U E_SU^*
\]
for a unitary \(U\) independent of \(n\).  The first congruence has
unitary equal to \(1\).  By the previously established tightness and
Lemma~\ref{coeff:analytic-calculus}, the families \(X'_n\) and
\((X'_n)^{-1}\), with both \(R\) and \(n\) varying, are tight.

Fix \(d,\varepsilon>0\).  If \(p=1_{(1+d,\infty)}(T)\) has
\(\tau(p)\ge\varepsilon\), then
\eqref{coeff:analytic-high} gives
\begin{equation}\label{coeff:analytic-gap}
 \|E_R\|_1+\|E_S\|_1
 \ge \tau\bigl(p\,\mathfrak M_n\delta_d(X'_n)\bigr),\qquad
 \delta_d(x)=\frac{dx}{(1+d+x)(1+x)}.
\end{equation}
Choose \(0<a<b<\infty\), independently of \(R,n\), such that
\[
 \tau(1-1_{[a,b]}(X'_n))<\varepsilon/2.
\]
On \([a,b]\), the function \(\delta_d\) has a positive minimum \(c\).
With \(q_n=1_{[a,b]}(X'_n)\), functional calculus gives
\[
 \tau(p\delta_d(X'_n))
 \ge c\,\tau(pq_n)
 \ge c\bigl(\tau(p)-\tau(1-q_n)\bigr)
 \ge c\varepsilon/2.
\]
These inequalities do not require \(p\) and \(q_n\) to commute.
They hold for every \(n\), so positivity of the scalar mean gives
the same lower bound in \eqref{coeff:analytic-gap}.  This contradicts
\eqref{coeff:analytic-approx-error} for sufficiently large \(R,S\).

For the low projection \(p=1_{(0,1-d)}(T)\), \(0<d<1\), use
\eqref{coeff:analytic-low} and the positive scalar function
\[
 \frac{1}{1+x}-\frac{1-d}{1-d+x}
       =\frac{dx}{(1+x)(1-d+x)}.
\]
The identical argument bounds its trace away from zero whenever
\(\tau(p)\ge\varepsilon\).  Thus \(T-1\to0\) in measure as \(R,S\to\infty\),
uniformly over these pairs.  The high and low estimates also imply
\(T^{-1}-1\to0\) in measure.  Since
\[
 P_S-P_R=P_R^{1/2}(T-1)P_R^{1/2},\qquad
 P_S^{-1}-P_R^{-1}
       =P_R^{-1/2}(T^{-1}-1)P_R^{-1/2},
\]
tightness and Lemma~\ref{coeff:analytic-calculus} show that both
\((P_R)\) and \((P_R^{-1})\) are Cauchy in measure.

Completeness gives limits \(P,V\).  Positivity is closed under
convergence in measure, and continuity of multiplication gives
\(PV=VP=1\).  Hence \(P\) is a positive affiliated unit and \(V=P^{-1}\).
Square-root continuity and the uniform product estimate show that
\[
 P_R^{-1/2}X_nP_R^{-1/2}
       -P^{-1/2}X_nP^{-1/2}\longrightarrow0
\]
in measure uniformly in \(n\).  The resolvent identity in
\eqref{coeff:analytic-F} gives uniform measure convergence of
\(F(P_R,X_n)\) to \(F(P,X_n)\).  These are contractions, so convergence
is uniform in \(L^1\).  Equations
\eqref{coeff:analytic-mean-L1} and
\eqref{coeff:analytic-approx-error} prove
\eqref{coeff:analytic-mean-equation}.

Finally, if \(P,Q\) are two exact solutions, their two-element families
and their inverses are tight.  Repeat the same congruence comparison
with zero errors.  Any nonzero high or low relative spectral
projection would give a strictly positive lower bound in
\eqref{coeff:analytic-gap}.  Therefore \(P^{-1/2}QP^{-1/2}=1\), proving
uniqueness.  The congruence formula then proves the asserted covariance:
a congruent solution solves the congruent equation, and it is the
unique such solution.
\end{proof}

\subsection{Proof of the support lemma}

\begin{proof}[Proof of Lemma~\ref{coeff:analytic-support}]
If \(\lambda=0\), the conclusion is immediate.  Assume otherwise.
The automorphism \(\alpha=\operatorname{Ad}g_D\) preserves \(L,S\),
their group von Neumann algebras, and their traces.  On affiliated
rows over \(L\), let
\[
 \mathcal T(\xi)=g\xi g_D^{-1}.
\]
This is well defined because \(h(g)=h(g_D)\).  It is right
\(\alpha\)-semilinear:
\(\mathcal T(\xi w)=\mathcal T(\xi)\alpha(w)\).

We first identify the column kernel with a projection in
\(M_r(\mathcal N(S))\).  Choose the positive affiliated unit
\[
 v=(1+\lambda\lambda^*)^{-1/2}\in\mathcal U(L),\qquad
 a=v\lambda.
\]
The row \(a\) is bounded.  Let
\[
 H=\mathbb E_S(a^*a)\in M_r(\mathcal N(S)),\qquad e=\supp H,
\]
where \(\mathbb E_S:\mathcal N(L)\to\mathcal N(S)\) is the
trace-preserving group conditional expectation, applied entrywise.
Concretely, if $P_S:\ell^2(L)\to\ell^2(S)$ is orthogonal projection,
then $\mathbb E_S(a)=P_Sa|_{\ell^2(S)}$. Compression commutes with
right $S$-translations and is positive and normal. Elements of
$\mathcal N(S)$ reduce $\ell^2(S)$, giving bimodularity; the vector
trace at $\delta_1$ gives trace preservation and faithfulness.
This is also the tracial expectation of
\cite[Theorem~9.1.2 and Remark~9.1.3]{AnantharamanPopa}.
Neither normality nor finite index of $S$ is required.
For a bounded \(S\)-column \(w\), compatible traces give
\[
 \tau((aw)^*aw)=\tau(w^*Hw).
\]
Faithfulness shows that \(aw=0\) if and only if \(ew=0\).
For an affiliated column, multiply on the right by increasing
spectral cuts for \(|w|\); the cut columns are bounded and converge
in measure to \(w\).  The same equivalence follows by measure
continuity.  Consequently
\begin{equation}\label{coeff:analytic-kernel}
 \ker(\lambda:\mathcal U(S)^r\to\mathcal U(L))
       =(1-e)\mathcal U(S)^r,\qquad \lambda e=\lambda.
\end{equation}
Here \(\lambda e=\lambda\) also follows directly from
\(a(1-e)=0\).  In particular \(e\ne0\).

The semilinear relation controls the transported row
\(g^n\lambda g_D^{-n}\). To apply the positive mean, we need
control of the coefficient maps themselves on the complement of the
column kernel. The following quantitative consequence of the kernel
description will transfer that control:
\begin{equation}\label{coeff:analytic-tight-reflection}
 \left.
 \begin{array}{c}
 w_j=ew_j,\\
 \{\lambda w_j\}\text{ is tight}
 \end{array}\right\}
 \quad\Longrightarrow\quad
 \{w_j\}\text{ is tight}.
\end{equation}
The matrices \(w_j\) here may have any fixed finite number of columns.
To prove it, suppose that
\(p=1_{(R,\infty)}(ww^*)\le e\) has \(\tau_r(p)\ge\varepsilon\).
Since \(\supp H=e\), choose \(\delta>0\) such that
\[
 \tau_r(e-1_{[\delta,\infty)}(H))<\varepsilon/2.
\]
Then, writing \(q=1_{[\delta,\infty)}(H)\),
\[
 \tau(apa^*)=\tau_r(pH)
 \ge\delta\,\tau_r(pq)
 \ge\delta\varepsilon/2=:c>0.
\]
Also \(0\le apa^*\le\|a\|^2\,1\).  Thus \(apa^*\) has spectral
mass at least \(c/(2\|a\|^2)\) above \(c/2\), using the canonical
trace \(\tau(1)=1\) in \(\mathcal N(L)\).
Positive order gives
\[
 aww^*a^*\ge Rapa^*.
\]
Monotonicity of spectral distribution under positive order therefore
forces the same positive amount of spectral mass above \(Rc/2\)
for \(aww^*a^*\).  This contradicts tightness of \(aw\) as
\(R\to\infty\).  Since \(aw=v(\lambda w)\), tightness of \(\lambda w\)
implies tightness of \(aw\) by
Lemma~\ref{coeff:analytic-calculus}.  This proves
\eqref{coeff:analytic-tight-reflection}.

Define the full invertible cocycle \(C_n\), \(n\in\Z\), by
\[
 C_0=1,\quad C_1=B,\quad
 C_{n+m}=C_n\alpha^n(C_m).
\]
Iterating \eqref{coeff:analytic-semilinear} gives
\[
 \mathcal T^n(\lambda)=g^n\lambda g_D^{-n}=\lambda C_n.
\]
Semilinearity and \eqref{coeff:analytic-kernel} imply
\[
 C_n\alpha^n\bigl((1-e)\mathcal U(S)^r\bigr)
                 =(1-e)\mathcal U(S)^r.
\]
Hence the induced maps on the quotient columns have representatives
\begin{equation}\label{coeff:analytic-partial-cocycle}
 b_n=eC_n\alpha^n(e):
       \alpha^n(e)\mathcal U(S)^r\longrightarrow e\mathcal U(S)^r.
\end{equation}
Their inverses are
\[
 d_n=\alpha^n(e)C_n^{-1}e=\alpha^n(b_{-n}),\qquad
 b_nd_n=e,\quad d_nb_n=\alpha^n(e).
\]
Compression does not spoil the cocycle law.  Indeed the discarded
summand \(\alpha^n(1-e)\) is carried by \(C_n\) into the kernel and
is killed by \(e\).  It follows that
\begin{equation}\label{coeff:analytic-compressed-law}
 b_{n+m}=b_n\alpha^n(b_m).
\end{equation}
Moreover,
\begin{equation}\label{coeff:analytic-compressed-orbit}
 \lambda b_n=\lambda C_n\alpha^n(e)
            =\mathcal T^n(\lambda).
\end{equation}

The spectral distributions of the rows on the right of
\eqref{coeff:analytic-compressed-orbit} are independent of \(n\):
their absolute squares are \(\alpha^n(\lambda^*\lambda)\).
Thus they form a tight family.  Since \(eb_n=b_n\),
\eqref{coeff:analytic-tight-reflection} shows that \(\{b_n\}\)
is tight.  Trace invariance of \(\alpha\) and the formula for \(d_n\)
show that \(\{d_n\}\) is tight as well.  In the fixed corner
\(eM_r(\mathcal U(S))e\), put
\[
 X_n=b_nb_n^*,\qquad X_n^{-1}=d_n^*d_n.
\]
Both families are tight, by the product estimate.

Apply Proposition~\ref{coeff:analytic-mean} in this corner to obtain
the unique positive affiliated unit \(P\) satisfying
\[
 \mathfrak M_nF(P,X_n)=\tfrac12 e.
\]
Equation~\eqref{coeff:analytic-compressed-law} gives
\[
 X_{n+1}=b_1\alpha(X_n)b_1^*.
\]
The automorphism \(\alpha\) commutes with the weak-* mean.
Its application moves the equation to the corner \(\alpha(e)\);
congruence by the affiliated invertible map
\(b_1:\alpha(e)\to e\) then gives an equation for
\(b_1\alpha(P)b_1^*\) with inputs \(X_{n+1}\).
Shift invariance and uniqueness yield
\begin{equation}\label{coeff:analytic-invariant-P}
 P=b_1\alpha(P)b_1^*.
\end{equation}

Set
\[
 Z=\lambda P^{1/2},\qquad
 u=P^{-1/2}b_1\alpha(P)^{1/2}.
\]
Inverses of \(P\) here and below are taken in the corner \(e\).
Equation~\eqref{coeff:analytic-invariant-P} gives
\(uu^*=e\) and \(u^*u=\alpha(e)\), so \(u\) is a bounded partial
unitary with entries in \(\mathcal N(S)\).  Equations
\eqref{coeff:analytic-compressed-orbit} and
\eqref{coeff:analytic-invariant-P} imply
\[
 gZg_D^{-1}=Zu,\qquad
 gZZ^*g^{-1}=Zuu^*Z^*=ZZ^*.
\]

Every conjugation-invariant positive affiliated operator in
\(\mathcal U(L)\) under this \(g\) is scalar.  To see this, a
nonidentity \(l\in L\) with finite orbit under conjugation by \(g\)
would commute with a nonzero power \(g^m\).
The centralizer of an infinite-order element in a hyperbolic group is
virtually cyclic; in the present torsion-free setting its
height-zero subgroup is trivial, since it contains \(g^m\) of nonzero
height.  Thus no such \(l\) exists.  Each spectral projection of an
invariant positive affiliated operator belongs to \(\mathcal N(L)\)
and has square-summable Fourier coefficients constant on these
conjugacy orbits.  All coefficients off the identity vanish.
Each spectral projection is therefore scalar, proving the assertion.
The resulting scalar is finite because the affiliated operator is
densely defined.

Applying this to \(ZZ^*\), we obtain \(ZZ^*=c1\), \(c\ge0\).
It follows that every entry \(Z_j\) is bounded, because
\(Z_jZ_j^*\le c1\).  Let \(\mathcal H\subseteq\ell^2(L)\) be the
closed right \(S\)-module generated by \(Z_1,\ldots,Z_r\).
The map
\[
 \ell^2(S)^r\longrightarrow\ell^2(L),\qquad
 (\xi_1,\ldots,\xi_r)\longmapsto\sum_{j=1}^r Z_j\xi_j
\]
is bounded and right \(S\)-equivariant, with range closure
\(\mathcal H\).  Its polar decomposition shows that the right
\(S\)-dimension of \(\mathcal H\) is at most \(r\)
\cite[Definitions~1.8 and~1.10, Theorem~1.9(6)]{Lueck2002}.
The unitary \(\mathcal T:\xi\mapsto g\xi g_D^{-1}\) on
\(\ell^2(L)\) preserves \(\mathcal H\): the formula
\(\mathcal T(Z)=Zu\) gives one inclusion, and
\(\mathcal T(Z)u^*=Z\) gives the other, using right semilinearity.

Decompose \(\ell^2(L)\) along the left cosets \(xS\), which are the
natural summands for its right \(S\)-action.  If \(\Pi\) is the
orthogonal projection onto \(\mathcal H\), its diagonal trace on
\(xS\) is
\[
 t_{xS}=\langle\Pi\delta_x,\delta_x\rangle\ge0.
\]
It is independent of the representative \(x\), and the module
dimension formula, valid also for infinite subgroup index, is
\[
 \sum_{xS\in L/S}t_{xS}=\dim_S\mathcal H\le r.
\]
The unitary \(\mathcal T\) permutes these cosets by
\[
 xS\longmapsto gxg_D^{-1}S.
\]
It commutes with \(\Pi\), so the diagonal traces are constant on
its permutation orbits.  Every infinite orbit therefore has trace
zero.  If \(t_{xS}=0\), positivity and right \(S\)-equivariance of
\(\Pi\) give \(\Pi\delta_{xs}=0\) for every \(s\in S\);
hence \(\mathcal H\) has zero Fourier component on that coset.
It follows that the coefficients of \(Z\) are supported on periodic
cosets.

A coset is periodic precisely when, for some \(m\ne0\),
\[
 g^mxg_D^{-m}S=xS.
\]
This implies \(x^{-1}g^mx\in D\).
Conversely, if \(x^{-1}g^mx\in D\), then
\(x^{-1}g^mxg_D^{-m}\) has height zero and lies in \(D\), hence lies
in \(S\); this proves the reverse implication.  Thus the periodic
cosets are exactly \(\mathcal P\).

Finally, \(\mathcal P\) is finite by the ambient quasiconvexity of
\(D\).  For \(xS\in\mathcal P\), left multiplication by a nonzero
power of \(g\) stabilizes the uniformly quasiconvex coset \(xD\).
Both endpoints \(g^+,g^-\) therefore lie in its limit set.
A geodesic joining these endpoints is contained in a uniformly
bounded neighborhood of \(xD\), with a bound depending only on
the hyperbolicity and quasiconvexity constants.  Fix a point on
this geodesic.  Every such \(xD\) meets one fixed finite ball
about that point, so there are finitely many such cosets \(xD\).
Since \(x\in L\),
\[
 xD\cap L=xS,
\]
and distinct \(xS\) in \(\mathcal P\) correspond to distinct \(xD\).
This proves finiteness.

Choose one representative \(x\in L\) for each coset in \(\mathcal P\).
Fourier projection, or equivalently the group conditional expectation,
now gives for each bounded entry
\[
 Z_j=\sum_{xS\in\mathcal P}x\,\mathbb E_S(x^{-1}Z_j),
 \qquad \mathbb E_S(x^{-1}Z_j)\in\mathcal N(S).
\]
This is an equality first in \(\ell^2(L)\), hence also in
\(\mathcal N(L)\).  Multiplying the row on the right by
\(P^{-1/2}\) recovers \(\lambda\), because \(\lambda e=\lambda\).
The coefficients become affiliated \(S\)-coefficients, proving
\eqref{coeff:analytic-support-conclusion}.
\end{proof}
